\section{Analysis: Why Better Token Fits Do Not Guarantee Better Agents}
\label{sec:analysis}

\paragraph{Gold-state localization versus generated-state control.}
Every locator score is computed on a verified trajectory prefix.  This is essential for causal comparability: hard-skill and no-skill distributions differ only in the skill.  During free generation, however, the prefix must first reproduce enough early structural decisions to remain on those states.  If it emits a different import, variable, workbook reference, or control-flow token, later training states become counterfactual.  This is the classical exposure-bias problem~\citep{bengio2015scheduled}; SpreadsheetBench amplifies it because semantic correctness is determined only after code execution.  Concentrated teacher effects can therefore coexist with weak rollout success.

\paragraph{The selected objective is local, while execution reward is structured.}
Equation~\ref{eq:combined} measures one-step distribution change, not a token's downstream causal effect.  A low-probability bracket or worksheet-selection call can determine whether hundreds of later tokens execute, while a high-JS lexical choice may have multiple semantically equivalent alternatives.  The multiplication $g_tj_t$ further suppresses positions where the hard skill changes the distribution but the cached gold token receives no positive gain.  Future locators should estimate downstream influence---for example, by teacher-forced suffix loss change, counterfactual execution, or on-policy visitation---without hand-designing fixed windows.

\paragraph{Why preservation helps but cannot solve the mismatch.}
The unselected KL prevents the prefix from arbitrarily degrading confident base behavior.  It is deliberately computed from a Top-64-plus-residual distribution rather than only the gold probability.  Yet it constrains only sampled gold states and only distributional drift.  It cannot guarantee syntactic invariants across a generated program, and it may oppose useful skill changes at unselected positions.  Its role is a trust region, not a proof of functional equivalence.

\paragraph{Why prefix-coordinate ``boosting'' mostly fails.}
Classic boosting adds weak learners in approximately independent function space~\citep{friedman2001greedy}.  Adjacent prefix vectors are not independent learners: through attention and every Transformer layer, updating one vector changes logits at nearly all subsequent tokens and changes the functional contribution of other vectors.  Consequently, PRCB-v1--v4 can reduce a block's monitored loss and still overwrite earlier behavior.  PRCB-v5's replay and line search reaches 18/40 validation tasks but only 96/280 test tasks.  PRCB-v6 creates independent full-length learners and combines their logits, but accepts only its first learner ($\alpha=.5$); the second stage obtains $\alpha=0$ even though 92.47\% of its selected core overlaps the previous core.  The residual remains large because the first learner has weak functional edge, not because the residual was computed without its prediction.  Distilling the ensemble back to one prefix does not improve validation; evaluating the undistilled ensemble also remains at 16/40.

\paragraph{What would materially change the optimization problem?}
The most promising next step is on-policy selective context distillation: sample student trajectories, query the hard-skill teacher on the exact student-visited states, select tokens subject to a capacity budget, and combine dense teacher KL with execution feedback.  This trades the clean causal interpretation of fixed gold contexts for coverage of the states encountered at inference.  A hybrid can retain verified gold anchors while adding student-state correction, related to on-policy context distillation~\citep{ye2026opcd} and generalized on-policy KD~\citep{agarwal2024gkd}.  This is a stronger intervention than another hand-chosen window or coordinate schedule.

\paragraph{Limitations.}
The current evidence is one model, one agent benchmark, one training seed, one teacher rollout per task, and a single final test evaluation.  The hard skill itself succeeds on only 61/80 training tasks, so the training set is success-conditioned.  Teacher generation uses a proprietary model, although trajectories are cached.  The main Combined configuration was chosen through validation comparisons, and the 280-task test was then used once; future changes must not reuse it for model selection.  Task-specific results are same-task oracles.  We report exact tests for paired outcomes but do not have multi-seed intervals.  These limitations prevent a broad state-of-the-art claim and define the experiments required before submission.
